% ED26-SRC-027-501, separate integration fragment for Remark3.11(iii).
% Integrator places this after the ordered-extension definition in section3.
\subsection{Global descent of the ordered marker extension}
\label{sec:p12-klop1987-ordered-marker-proof}

Remark3.11(iii), printed p.22 of the 1987 report, changes both the
objects and the marker operations: sibling order matters, each operator
has a fixed arity, pushing changes just one argument, and simplifying
the left argument is added \parencite{Klop1987CWITutorial}. A proof of
the earlier commutative game does not establish this enlarged game's
well-foundedness. The following supplies that missing argument for
the actual fixed-arity operations. Its dependency is the fully proved
left-to-right LPO of Definitions~\ref{def:p12-ro-data} and
\ref{def:p12-ro-lpo}, with Theorems~\ref{thm:p12-ro-lpo-structural},
\ref{thm:p12-ro-finite-tree} and \ref{thm:p12-ro-lpo-termination}.
The primary Kamin--L\'evy note distinguishes bounded lexicographic
comparison from an unrestricted tuple comparison: the source must also
dominate every target argument
\parencite[handwritten pp.1--8]{KaminLevy1980OrderedMarkers}.
Here that domination is established even when different marked copies
have undergone different reductions; it is not assumed from their names.

\begin{hbdefinition}{def:p12-klop1987-om-game}{Fixed-arity ordered trees and all five root operations}
Let \(\Sigma\subseteq\mathbb N\), with one prescribed finite arity
\(a(n)\) for each \(n\in\Sigma\). Plain trees are finite, nonempty
ground terms over this ranked signature. Marked trees use \(n\) or
\(n^*\) at a node of arity \(a(n)\), without permuting children.
For a tree with plain root \(u=n(\vec v)\), let
\(u^*=n^*(\vec v)\). As in source Notation3.2, this operation is
defined at a plain root; it does not add a second star or silently
interpret a star on an already starred root as idempotent.
Write \(k=a(n)\). The root operations are
\begin{align*}
 \mathrm{PUT}:\quad &n(\vec u)\longrightarrow n^*(\vec u),\\
 \mathrm{SELECT}:\quad &n^*(u_1,\ldots,u_k)\longrightarrow u_i
       &&(1\leq i\leq k),\\
 \mathrm{COPY}:\quad &n^*(\vec u)\longrightarrow
       m(\underbrace{n^*(\vec u),\ldots,n^*(\vec u)}_{a(m)})
       &&(m\in\Sigma,\ m<n),\\
 \mathrm{PUSH}:\quad &n^*(u_1,\ldots,u_i,\ldots,u_k)
       \longrightarrow n(u_1,\ldots,u_i^*,\ldots,u_k),\\
 \mathrm{LEX}:\quad &n^*(u_1,\ldots,u_k)\longrightarrow
       n(u_1^*,\underbrace{n^*(u_1,\ldots,u_k),\ldots,
                      n^*(u_1,\ldots,u_k)}_{k-1})
       &&(k\geq1).
\end{align*}
PUSH requires \(1\leq i\leq k\) and a plain-root \(u_i\);
LEX requires a plain-root \(u_1\), exactly so that the displayed
root-mark operations are defined. Other descendant marks are permitted.
COPY has zero copies exactly when \(a(m)=0\); its copy count is not
free. PUSH neither copies nor deletes its chosen argument. LEX at
arity one has only \(u_1^*\) on its right-hand argument list.

Denote the source's closure under plain-root contexts by
\(\to_{\rm lex}\), and the expressly larger closure under all
finite marked or plain contexts by \(\rightsquigarrow_{\rm lex}\).
The source context clause puts plain labels on the path above the
hole; its fixed sibling subtrees may themselves contain marks.
Every source step is an auxiliary step. The theorem below proves
descent for this larger game, hence for the stated source game.
\end{hbdefinition}

\begin{hblemma}{thm:p12-klop1987-om-bound}{proved=annex:p12-klop1987-om-bound}{One finite ranked signature for an entire run}
If the largest numerical label in an initial marked or plain tree is
\(M\), every descendant under the auxiliary game uses only
\(\Sigma_M=\{n\in\Sigma:n\leq M\}\). This is one fixed finite
ranked signature, with at most \(M+1\) symbols, throughout the run.
The assertion bounds neither the number of nodes nor depth nor the
number of copies or steps. In particular it is not a finite-depth
special case of the ordered extension.
\end{hblemma}

\begin{hbdefinition}{def:p12-klop1987-om-interpretation}{Nonempty weak and strict LPO lower sets}
Fix \(M\) for a run and add a fresh nullary symbol \(\bot\) below
every member of \(\Sigma_M\). It is only a semantic symbol: the
game neither contains nor marks this bottom symbol. Let \(U_M\) be
the ground terms over \(\Sigma_M\cup\{\bot\}\), with total
precedence inherited from the numerical labels and bottom last, and
left-to-right LPO. Write \(>\) for this LPO and \(\geq\) for
equality or \(>\). Define
\[
 D(v)=\{w\in U_M:v\geq w\},\qquad
 D^-(v)=\{w\in U_M:v>w\}.
\]
For a marked game tree define its interpretation recursively by
\begin{align*}
 J(n(u_1,\ldots,u_k))
  &=\bigcup_{b_i\in J(u_i)}D(n(b_1,\ldots,b_k)),\\
 J(n^*(u_1,\ldots,u_k))
  &=\bigcup_{b_i\in J(u_i)}D^-(n(b_1,\ldots,b_k)).
\end{align*}
The product over an empty argument list has one empty tuple.
A lower set means downward closed for \(\geq\). Lower sets here
may be infinite and need not have greatest elements.
\end{hbdefinition}

\begin{hblemma}{thm:p12-klop1987-om-sets}{proved=annex:p12-klop1987-om-sets}{Well-defined interpretations, context monotonicity and finite witnesses}
The interpretation of every finite marked tree is a nonempty lower
set containing \(\bot\). Both constructor interpretations are
monotone under inclusion of each argument set. Consequently replacing
a subtree by one with a smaller interpretation cannot enlarge the
interpretation in any marked or plain context.
Every finite nonempty family of members of one such argument set has
a greatest member belonging to that set, since ground LPO is total.
Thus finitely many argument-choice witnesses can be dominated
coordinatewise within the same argument sets. This uses maxima of
finite families, not a greatest member of an entire lower set.
\end{hblemma}

The added bottom repairs an actual failure of the naive interpretation.
Without it, over a signature with nullary \(c\) and binary \(n>c\),
the least term \(c\) has \(D^-(c)=\varnothing\). The naive
interpretations would then give \(J(n^*(c^*,c))=\varnothing\)
but \(J(c)=\{c\}\), although SELECT of the second child is
\(n^*(c^*,c)\to c\). Thus downward closure alone, allowing empty
argument sets, would not justify operation monotonicity.

\begin{hblemma}{thm:p12-klop1987-om-operations}{proved=annex:p12-klop1987-om-operations}{Every operation decreases the lower-set interpretation}
For each of PUT, SELECT, COPY, PUSH and LEX, and hence for every
\(u\rightsquigarrow_{\rm lex}v\), one has \(J(v)\subseteq J(u)\).
COPY includes its nullary case. The LEX assertion covers every fixed
arity \(k\geq1\), with independent interpretation choices in its
\(k-1\) marked copies. No equality between those choices is required.
\end{hblemma}

\begin{hblemma}{thm:p12-klop1987-om-first-mark}{proved=annex:p12-klop1987-om-first-mark}{The first mark gives strict endpoint descent}
For an entirely plain tree \(t\), \(J(t)=D(t)\). For its root
mark \(t^*\), \(J(t^*)=D^-(t)\). If a first PUT is performed
at any position of an entirely plain tree \(t\), producing \(u\),
then \(J(u)\subseteq D^-(t)\). Therefore every positive finite
auxiliary reduction between plain trees satisfies
\[
 t\rightsquigarrow_{\rm lex}^{+}s\quad\Longrightarrow\quad t>s.
\]
The individual auxiliary steps need only preserve inclusion; the
first mark supplies the strictness retained at the plain endpoint.
\end{hblemma}

\begin{hbtheorem}{thm:p12-klop1987-om-well-founded}{proved=annex:p12-klop1987-om-well-founded}{Global well-foundedness of the ordered fixed-arity marker order}
For every ranked signature \(\Sigma\subseteq\mathbb N\), without
one common global label, arity, size or depth bound, the positive
auxiliary relation restricted to pairs of plain trees is transitive,
irreflexive and well-founded. The same holds for the source relation.
In particular no infinite sequence of plain trees can be connected by
successive positive finite reductions using the ordered LEX extension.
This is the missing global conclusion for Remark3.11(iii), not just
an orientation of the report's displayed Ackermann pair.
\end{hbtheorem}

This conclusion concerns plain endpoints, not termination of the whole
marked universe. For a binary \(n\), a plain-root \(c\) and any
\(u\), the auxiliary game has the marked cycle
\[
 n^*(c,u)\longrightarrow n(c^*,n^*(c,u))
 \longrightarrow n^*(c^*,n^*(c,u))\longrightarrow n^*(c,u).
\]
The steps are LEX, PUT and SELECT of the second child. Its endpoints
are marked, so it contradicts neither the theorem nor the strict
plain-endpoint interpretation.

\begin{hbtheorem}{thm:p12-klop1987-om-applications}{proved=annex:p12-klop1987-om-applications}{All source Ackermann rules and the ordered associativity example terminate}
The report's three rules
\[
 A(0,x)\to S(x),\qquad A(S(x),0)\to A(x,S(0)),\qquad
 A(S(x),S(y))\to A(x,A(S(x),y))
\]
have positive ordered-marker simulations with fixed arities
\(a(A)=2\), \(a(S)=1\), \(a(0)=0\) and labels
\(A=2>S=1>0\). They terminate on all finite ordinary terms,
over \(\{A,S,0\}\) with variables, not just closed numerical
inputs. The rule
\((x\mathbin{\cdot}y)\mathbin{\cdot}z\to
x\mathbin{\cdot}(y\mathbin{\cdot}z)\) likewise has a positive
ordered-marker simulation with one fixed binary product operator,
and terminates. The annex supplies the actual paths and transports
them through arbitrary rule instances and contexts.
\end{hbtheorem}

\begin{proofboundary}{def:p12-klop1987-om-standing}{What this proof establishes and what remains a different method}
The lower-set argument is an editorial completion of the 1987 marker
presentation, using the separately proved LPO dependency. It is not
claimed to be printed in full in the tutorial, in the undated primary
Kamin--L\'evy scan, or in the distinct 1992 chapter. Its finite-signature
step is justified by the marker run's numerical-label invariant, not
imposed as a restriction on the theorem's infinite ambient signature.
It neither reidentifies ordered fixed-arity trees with the earlier
commutative unranked trees, nor imports a stronger embedding convention
as the report's own definition. Puel's extension for an embedded left
side and the full norm-indexed process-algebra development remain
separate methods and obligations.
\end{proofboundary}
